La figure ci-dessous donne deux dispositifs~(1) et~(2) permettant d'étudier la
propagation d'une onde à la surface de l'eau et la propagation du son dans
l'air.

%\begin{figure}[H]
%    \centering
%    \begin{subfigure}{0.45\textwidth}
%        \centering
%        \includegraphics[width=\textwitdh]{g8136.png}
%        \caption{}
%        \label{fig:}
%    \end{subfigure}
%    \begin{subfigure}{0.45\textwidth}
%        \centering
%        \includegraphics[width=\textwitdh]{g8137.png}
%        \caption{}
%        \label{fig:}
%    \end{subfigure}
%    \caption{}
%    \label{fig:}
%\end{figure}

\begin{figure}[H]
    \centering
    \includegraphics[width=0.7\textwidth]{2025-11-01_012313-}
\end{figure}

\begin{enumerate}
    \item Quelle est la nature de l'onde mécanique produite respectivement par
          les sources de ces deux dispositifs~?

    \item Dans le dispositif~(1), un vibreur produit une onde progressive
          sinusoïdale de fréquence $N_{1}$. Une étude expérimentale a permis
          d'obtenir le document~(a) représentant l'élongation d'un point $M$ de
          la surface de l'eau en fonction du temps et le document~(b)
          représentant l'aspect de la surface de l'eau à un instant donné.
          \begin{figure}[H]
              \centering
              \includegraphics[width=0.8\textwidth]{2025_10_31_756318f153383b902639g-1}
          \end{figure}
          \begin{enumerate}
              \item Lequel des deux documents~(a) et~(b) montre une
                    périodicité spatiale~?
              \item Déterminer la fréquence $N_{1}$ de l'onde.
              \item Calculer la célérité $v_{1}$ de propagation de l'onde à la
                    surface de l'eau.
              \item Recopier sur votre copie le numéro de la question et écrire
                    la lettre correspondante à la proposition vraie.
              
                    L'élongation du point $M$ en fonction de l'élongation de la
                    source $S$ s'écrit~:
                    \begin{table}[H]
                        \flushright
                        \setlength{\arrayrulewidth}{1pt}
                        \renewcommand{\arraystretch}{1.5}
                        \begin{tabularx}{0.95\textwidth}{|p{3mm}|X|p{3mm}|X|p{3mm}|X|p{3mm}|X|}
                            \hline
                            $A$ & $y_{M}(t)=y_{S}(t+0,1)$ &
                            $B$ & $y_{M}(t)=y_{S}(t+0,05)$ &
                            $C$ & $y_{M}(t)=y_{S}(t-0,1)$ &
                            $D$ & $y_{M}(t)=y_{S}(t-0,05)$ \\
                            \hline
                        \end{tabularx}
                    \end{table}
          \end{enumerate}
    \item On interpose à la surface de l'eau un obstacle muni d'une ouverture de
          diamètre $L=\qty{8}{\cm}$. L'onde produite à la surface de l'eau par
          la source se propage après avoir traversé l'ouverture.
          \begin{enumerate}
              \item Quel phénomène peut-on observer lorsque l'onde traverse
                    l'ouverture~? Justifier.
              \item Déduire la longueur d'onde $\lambda_{2}$ et la célérité de
                    propagation $v_{2}$ de l'onde au-delà de l'ouverture.
          \end{enumerate}
    \item Le haut-parleur du dispositif~(2), émet des ondes sonores de fréquence
          $N_{2}=\qty{10}{\kHz}$.
          \begin{enumerate}
              \item Les ondes sonores produites peuvent-elles se propager dans
                    le vide? Justifier.
              \item Les ondes sont captées par deux microphones $M_{1}$
                    et $M_{2}$ qui occupent la même position. Les
                    courbes visualisées sur l'écran de l'oscilloscope
                    apparaissent en phase.
              
                    Lorsqu'on déplace $M_{2}$ par rapport à
                    $M_{1}$ d'une distance $d=\qty{34}{\cm}$, les deux courbes
                    observées à l'oscilloscope apparaissent à nouveau en phase
                    pour la dixième~(10) fois. Déduire la célérité de
                    propagation du son dans l'air.
          \end{enumerate}
\end{enumerate}

